
\sourcenote{Genre: alias rule (tier 3) --- a respelling template. An alias
owes the book five artefacts: the verified 1$\to$1 unpack, the filled-in
application of the target rule, the payload-transport table, the exact
algebraic identity that makes the respelling correct, and the soundness
corollary through the target rule's lemma. All five follow --- and this
alias is the tier taxonomy's \emph{canonical} case, the one the
tier was named for: the unpack is the identity.}

\begin{brecipe}[the alias template]\label{rec:pb-template}
An alias rule $A$ is warranted by a verified function
$\mathrm{unpack} : \mathrm{Step}_A \to \mathrm{Step}_T$ into the payload
of a \emph{single} target node $T$ (core or admissible), such that the
$A$-node and its image conclude the same sequent from the same premises:
one node $\to$ one node, payload transported, nothing derived. Soundness
is then a corollary, $\mathrm{sound}_A = \mathrm{sound}_T \circ
\mathrm{unpack}$. The wire format may carry $A$ as a compact register of
its own (this rule: an affine coefficient-row spelling, tag pending
allocation) or never carry it at all (producers elaborating before
emission, as the \rulename{SmallGain} section records); either way the
checker's obligations are the target's, reached through the unpack.
\end{brecipe}

\begin{brecipe}[the verified unpack --- the identity transport]
\label{rec:pb-unpack}
\rulename{PBound}'s payload carries components, combination term,
conclusion term, and the bookkeeping fields \emph{verbatim} as the
target's payload does, including the conclusion-form selector. The
unpack is therefore not a computation but a re-labelling: the verified
object is the direct specialization
\[
\mathrm{sound}_{\rulename{PBound}}
\;=\;
\mathrm{sound}_{\rulename{Syzygy}}
\ \text{\upshape applied to the same fields, re-assembled as the
target's step}
\]
--- in the sources, the body of \texttt{sound\_pBound} \emph{is} one
application of \texttt{sound\_syzygy} to the step built from
\rulename{PBound}'s own fields. Where \rulename{SmallGain}'s unpack
computes a weighting, this one computes nothing; that is the tier-3
boundary case made flesh, and it is why the alias relation here is
$1 \to 1$ at the level of \emph{readings} too
(Recipe~\ref{rec:pb-transport}).
\end{brecipe}

\begin{brecipe}[the filled-in application]\label{rec:pb-filled}
Under the identity transport the \rulename{PBound} node \emph{is} the
following application of the target rule (compare the rule display in
the \rulename{Syzygy} section, field for field):
\begin{mathpar}
\inferrule*[Right=\rulename{Syzygy}]{
  \bigl\{\Phi_i \vdash e_i \leq 0\bigr\}_{i=1}^{n} \\
  \bigl\{g_k = 0\bigr\}_k \ \text{sub-conjuncts of } \Phi \\\\
  \lambda_i \in \Q_{\geq 0} \\
  \nu_k \in \Q \\
  c \;\equiv\; \textstyle\sum_i \lambda_i e_i + \sum_k \nu_k g_k
  \ \text{(exact canonical polynomial)}
}{
  \Phi \vdash \Psi
}
\end{mathpar}
where $\Psi$ is whichever reading of $c \leq 0$ the shared selector
declares: \emph{(Z$\leq$)} $c \leq 0$ itself; \emph{(G)} $e \geq 0$
where $c = -e$; \emph{(L$\leq$)} $l \leq r$ where $c = l - r$.
Discharging the target's obligations at this instance is immediate ---
they are \emph{the same obligations on the same data}: entry decoding,
nonnegative weights, entry-sourcing (component conclusions and premise
equality sub-conjuncts, both polarities on equalities), the exact fold
identity, premise containment, conclusion shape/finiteness, and the
positive-size wire conditions on premise, conclusion, and component
conclusions. \rulename{PBound}'s catalogued statement restricts the
$e_i$ to \emph{affine} entries; nothing above consults that restriction
(Remark~\ref{rem:pb-affine}).
\end{brecipe}

\begin{brecipe}[payload transport --- field for field]
\label{rec:pb-transport}
\begin{center}\small
\begin{tabular}{@{}ll@{}}
\textbf{\rulename{PBound} payload} & \textbf{target (\rulename{Syzygy}) payload}\\
\midrule
\texttt{components} & \texttt{components} (verbatim)\\
\texttt{combination} & \texttt{combination} (verbatim)\\
\texttt{conclusion} & \texttt{conclusion} (verbatim)\\
\texttt{denBoundLog2} & \texttt{denBoundLog2} (verbatim; witness-side claim, zero proof-theoretic weight)\\
\texttt{preContains} & \texttt{preContains} (verbatim)\\
\texttt{conclShape, conclFin} & \texttt{conclShape, conclFin} (verbatim)\\
\texttt{wiresPos, compPostWiresPos} & \texttt{wiresPos, compPostWiresPos} (verbatim)\\
\texttt{valid} & \texttt{valid} (the same exact-combination proposition)\\
\texttt{postForm} & \texttt{postForm} (the same three-reading selector)\\
\end{tabular}
\end{center}
Every row is the identity --- including the last: since the target rule
itself carries the three-reading conclusion selector, the alias
inherits all readings with no residue, and the transport is $1 \to 1$
at reading level, not merely at node level.
\end{brecipe}

\begin{bprop}[the exactness identity]\label{prop:pb-exact}
The respelling is exact definitionally: the unpacked node's combination
identity is \emph{the same} exact canonical-polynomial identity
$c \equiv \sum_i \lambda_i e_i + \sum_k \nu_k g_k$, unchanged. What
remains to make the three readings correct is the fact that they are
one exact $\Q$ identity re-presented under a sign flip --- per entry,
$-e(\sigma) \leq 0 \iff e(\sigma) \geq 0$ and
$l(\sigma) - r(\sigma) \leq 0 \iff l(\sigma) \leq r(\sigma)$ --- which
is the target rule's reading-transport lemma, proven once in the
\rulename{Syzygy} section and \emph{not} re-justified here.
\end{bprop}
\begin{proof}
The first claim is reflexivity of the transported field. The second is
the cited lemma (the ordered-field axioms negate-and-flip and
add-to-both-sides, packaged as the ge-bridge ring lemma); the alias
contributes no polarity argument of its own, which is what keeps the
$1 \to 1$ relation clean.
\end{proof}

\begin{bprop}[soundness corollary]\label{prop:pb-sound}
If every component sequent is semantically valid, then under the
payload's checked hypotheses the sequent $\Phi \vdash \Psi$ is
semantically valid, for every reading $\Psi$ the selector admits.
\end{bprop}
\begin{proof}
\texttt{sound\_pBound} $=$ \texttt{sound\_syzygy} at the transported
(identical) payload data, fed the same component validities. There is
no other content --- the alias genre working as intended.
\end{proof}

\begin{bremark}[why the rule exists at all]\label{rem:pb-why}
Semantically the alias adds nothing to its target; its existence is
pragmatic, twice over. \emph{Register}: affine entries admit a compact
coefficient-row wire spelling, and bound nodes are the bulk commodity
of produced certificates --- the compact register is worth a tag.
\emph{Readings at the seams}: the (G) reading makes a nonnegativity
premise (as \rulename{ReLU-Pos} demands, $h \geq 0$) reachable
directly from bound facts --- the once-contemplated standalone
$\geq$-bridge rule dissolved into the conclusion selector instead ---
and the (L$\leq$) reading hands interval-shaped consumers
(\rulename{LiRPA}'s endpoint premises) their literal atom with no
\rulename{ArithAtom} re-arrangement node in between. Producers get the
convenient conclusion shape; the kernel keeps exactly one combination
argument.
\end{bremark}

\begin{bremark}[the affine restriction is register, not soundness]
\label{rem:pb-affine}
The catalogued statement's ``affine entries'' names the wire register
the rule is priced for, and is not a checked side condition: the
combination replay is degree-blind, and the soundness argument
(Proposition~\ref{prop:pb-sound}) never consults linearity. A
higher-degree entry simply belongs to the target rule's own spelling.
Publishing the restriction in the rule head, while the checker enforces
only what soundness needs, mirrors the \rulename{Farkas} page's
linearity note --- the register tells producers what the format prices;
the checker tells them what truth requires.
\end{bremark}

\begin{anchors}
Recipe \ref{rec:pb-unpack} unpack (identity transport) & \texttt{sound\_pBound}'s body: one application of \texttt{sound\_syzygy} to the step assembled from \rulename{PBound}'s fields (\texttt{PBoundData} $=$ the target's fields verbatim)\\
Recipe \ref{rec:pb-filled} obligations & \texttt{syzygyCombinationValid} and the target's bookkeeping fields, unchanged\\
Recipe \ref{rec:pb-transport} reading selector & \texttt{syzygyConclusionForm} (shared, not duplicated)\\
Proposition \ref{prop:pb-exact} reading transport & \texttt{holds\_of\_syzygyConclusionForm} (proven in the \rulename{Syzygy} section)\\
Proposition \ref{prop:pb-sound} & \texttt{sound\_pBound} (a \texttt{sound\_syzygy} specialization)\\
wire placement & \texttt{RuleTag.pBound}; tag allocation deferred --- the emit path refuses the rule until the payload-window allocation\\
\end{anchors}
